% !TEX program = xelatex
% 공통수학2 · Ⅰ. 도형의 방정식 · 2. 원의 방정식 · 01 원의 방정식 · 2/2차시
% 학생용 활동지 (답안 없음)
\documentclass[11pt,a4paper]{article}
\usepackage{kotex}
\usepackage{iftex}
\ifPDFTeX\else
  \setmainhangulfont{Noto Serif CJK KR}
  \setsanshangulfont{Noto Sans CJK KR}
\fi
\usepackage[margin=2.1cm]{geometry}
\usepackage{parskip}
\usepackage{enumitem}
\usepackage{array}
\usepackage{tabularx}
\usepackage{xcolor}
\usepackage{amssymb}
\usepackage{tikz}
\usepackage[most]{tcolorbox}
\usepackage{fancyhdr}
\usepackage{titlesec}

\definecolor{goldc}{HTML}{B07C22}
\definecolor{goldstrong}{HTML}{8A5F16}
\definecolor{indigoc}{HTML}{2B3B57}
\definecolor{indigosoft}{HTML}{6E82A6}
\definecolor{linec}{HTML}{D6DACB}
\definecolor{surface2}{HTML}{E4E8DC}
\definecolor{writeline}{HTML}{C7CCBA}
\definecolor{inksoft}{HTML}{52604F}

\titleformat{\section}{\normalfont\sffamily\bfseries\large\color{indigoc}}{\thesection}{0.6em}{}
\titlespacing{\section}{0pt}{14pt}{6pt}
\newtcolorbox{goalbox}{enhanced, breakable, colback=white, colframe=goldc, boxrule=0.5pt, leftrule=3pt, arc=2pt, left=10pt, right=10pt, top=8pt, bottom=8pt}
\newtcolorbox{sheetbox}{enhanced, breakable, colback=white, colframe=linec, boxrule=0.6pt, arc=3pt, left=11pt, right=11pt, top=9pt, bottom=9pt}
\newcommand{\writespace}[1]{%
  \par\nobreak\vspace{3pt}
  \foreach \n in {1,...,#1} {%
    \noindent\textcolor{writeline}{\rule{\linewidth}{0.4pt}}\par\vspace{0.62cm}%
  }
}
\newcommand{\fillblank}[1]{\makebox[#1]{\hrulefill}}
\pagestyle{fancy}\fancyhf{}\renewcommand{\headrulewidth}{0pt}
\fancyfoot[L]{\footnotesize\color{inksoft} 공통수학2 · 2026학년도 2학기 · 지평선고등학교}
\fancyfoot[R]{\footnotesize\color{inksoft} 다음 시간 --- 02 원과 직선의 위치 관계(1)}

\begin{document}

\noindent
\begin{minipage}[b]{0.62\linewidth}
  {\sffamily\footnotesize\color{indigosoft} 공통수학2 · Ⅰ. 도형의 방정식 · 2. 원의 방정식}\\[2pt]
  {\sffamily\bfseries\Large 01 원의 방정식 \textperiodcentered\ 10차시}
\end{minipage}\hfill
\begin{minipage}[b]{0.36\linewidth}
  \raggedleft\small\color{inksoft}
  반\ \fillblank{1.6cm} \quad 번호\ \fillblank{1.6cm}\\[4pt]
  이름\ \fillblank{3.6cm}
\end{minipage}

\vspace{4pt}
\noindent{\color{goldc}\rule{\linewidth}{2.2pt}}
\vspace{10pt}

\begin{goalbox}
\textbf{오늘의 목표}\\
세 점을 지나는 원의 방정식을 구하고, 이차방정식 $x^2+y^2+Ax+By+C=0$이 나타내는 도형(원)을 표준형으로 바꿀 수 있다.
\vspace{4pt}
\small\color{inksoft}
$\square$ 자신 있음 \quad $\square$ 조금 헷갈림 \quad $\square$ 처음 봄 \hfill (수업 전 나의 상태에 표시)
\end{goalbox}

\section{생각 열기 --- 세 점으로 원 전체 알아내기}

\begin{sheetbox}
세 점 A$(1,5)$, B$(2,4)$, C$(-1,1)$이 한 원 위에 있다.

\begin{enumerate}[leftmargin=1.4em, itemsep=2pt]
  \item 원의 중심을 P$(a,b)$라 하면, 중심에서 세 점까지의 거리는 어떤 관계여야 할까?
  \item 이 관계를 식으로 나타내면 몇 개의 방정식을 얻을 수 있을까?
\end{enumerate}
\writespace{3}
\end{sheetbox}

\section{개념 정리 --- 세 점을 지나는 원 \& 일반형}

\begin{sheetbox}
\textbf{세 점을 지나는 원} --- 중심을 P$(a,b)$라 하면 $\overline{\text{PA}}=\overline{\text{PB}}=\overline{\text{PC}}$이므로
\[
\overline{\text{PA}}^2 = \fillblank{2cm} \,,\qquad \overline{\text{PB}}^2 = \fillblank{2cm}
\]
두 식을 연립하여 $a,b$를 구한다.

\vspace{10pt}
\textbf{일반형} --- 원의 방정식 $(x-a)^2+(y-b)^2=r^2$을 전개하여 정리하면
\[
x^2+y^2+Ax+By+C=0
\]
꼴로 나타낼 수 있다. 이 식을 완전제곱식으로 변형하면 중심 $\left(-\dfrac{A}{2},-\dfrac{B}{2}\right)$, 반지름 $\dfrac{\sqrt{A^2+B^2-4C}}{2}$인 원이 된다. 단, 이 식이 원이 되려면
\[
A^2+B^2-4C \ \fillblank{1.6cm}\ 0
\]
이어야 한다.
\end{sheetbox}

\section{연습 문제}

\begin{sheetbox}
\textbf{1.} 세 점 A$(-2,6)$, B$(-5,-3)$, C$(2,-2)$를 지나는 원의 방정식을 구하시오.
\writespace{5}

\vspace{6pt}
\textbf{2.} 다음 방정식이 나타내는 도형을 말하시오.\\
\hspace*{1.6em}⑴ $x^2+y^2+4x=0$ \qquad ⑵ $x^2+y^2-10x-8y-8=0$
\writespace{4}

\vspace{6pt}
\textbf{3.} 방정식 $x^2+y^2-2x+6y+k=0$이 나타내는 도형이 원이 되도록 하는 실수 $k$의 값의 범위를 구하시오.
\writespace{3}
\end{sheetbox}

\section{사고의 가시화 --- Connect · Extend · Challenge}

\noindent{\footnotesize\color{inksoft} `원의 방정식' 전체(9$\sim$10차시)를 떠올리며 아래 세 칸을 채워 보자.}\\[6pt]
\begin{tabularx}{\linewidth}{@{}XXX@{}}
\textbf{\footnotesize\color{indigoc} Connect --- 무엇과 연결되나?} &
\textbf{\footnotesize\color{indigoc} Extend --- 어디까지 확장됐나?} &
\textbf{\footnotesize\color{indigoc} Challenge --- 아직 궁금한 것은?} \\[4pt]
\end{tabularx}
\noindent
\begin{minipage}[t]{0.32\linewidth}\writespace{3}\end{minipage}\hfill
\begin{minipage}[t]{0.32\linewidth}\writespace{3}\end{minipage}\hfill
\begin{minipage}[t]{0.32\linewidth}\writespace{3}\end{minipage}

\section{마무리 성찰}

\begin{sheetbox}
\begin{minipage}[t]{0.48\linewidth}
{\footnotesize\color{inksoft} `원의 방정식' 소단원을 마치며 한 문장으로}
\writespace{2}
\end{minipage}\hfill
\begin{minipage}[t]{0.48\linewidth}
{\footnotesize\color{inksoft} 감사 일기 / 유념 한 줄}
\writespace{2}
\end{minipage}
\end{sheetbox}

\end{document}
